\documentclass[11pt]{article}

\usepackage[margin=1.8cm]{geometry}
\usepackage{amsmath}
\usepackage{amssymb}
\usepackage{enumitem}
\usepackage{titlesec}
\usepackage{xcolor}
\usepackage{fancyhdr}
\usepackage{multicol}

\titleformat{\section}{\large\bfseries}{\thesection}{0.6em}{}
\titlespacing{\section}{0pt}{1.2em}{0.6em}

\pagestyle{fancy}
\fancyhf{}
\fancyhead[L]{Solving Exponential Equations}
\fancyhead[R]{Name: \makebox[4cm]{\hrulefill}}
\fancyfoot[C]{\thepage}
\renewcommand{\headrulewidth}{0.4pt}

\newcommand{\worksheettitle}{Solving Equations with Exponentials}
\newcommand{\qspace}{\vspace{0.65cm}}

\setlist[enumerate,1]{leftmargin=1.4em, itemsep=0.3em, topsep=0.15em}
\setlist[enumerate,2]{leftmargin=1.2em, itemsep=0.25em, topsep=0.15em}

\setlength{\columnsep}{1.1cm}
\setlength{\columnseprule}{0.4pt}

\newenvironment{qtwocol}{\begin{multicols}{2}\begin{enumerate}[resume]}{\end{enumerate}\end{multicols}}
\newenvironment{qtwocolstart}{\begin{multicols}{2}\begin{enumerate}}{\end{enumerate}\end{multicols}}

\begin{document}

\begin{center}
  {\Large\bfseries \worksheettitle}\\[0.3em]
  {\normalsize Fluency \(\cdot\) Problem Solving \(\cdot\) Enrichment}
\end{center}

\vspace{0.3em}
\noindent\textit{Big idea: an exponential equation can be solved without calculus or logarithms as soon as both sides are written as a power of the \textbf{same base}. Once that is true, the exponents themselves must be equal, turning the problem into an ordinary equation.}

\vspace{0.4em}

\section{Fluency}

\subsection*{Multiplying exponentials ($a^m \times a^n = a^{m+n}$)}
Simplify each expression.

\begin{qtwocolstart}
  \item $2^3 \times 2^4$
  \qspace
  \item $x^5 \times x^2$
  \qspace
  \item $5^2 \times 5^4 \times 5$
  \qspace
  \item $4^2 \times 4^5$
  \qspace
  \item $3 \times 3^4 \times 3^2$
  \qspace
  \item $x^3 \times x \times x^4$
  \qspace
\end{qtwocolstart}

\subsection*{Dividing exponentials ($a^m \div a^n = a^{m-n}$)}
Simplify each expression.

\begin{qtwocol}
  \item $\dfrac{5^7}{5^3}$
  \qspace
  \item $\dfrac{x^9}{x^4}$
  \qspace
  \item $\dfrac{3^8 \times 3^2}{3^5}$
  \qspace
  \item $\dfrac{6^{10}}{6^{6}}$
  \qspace
  \item $\dfrac{7^{12}}{7^{5}}$
  \qspace
  \item $\dfrac{x^{11}}{x^{3}}$
  \qspace
\end{qtwocol}

\subsection*{The power law ($(a^m)^n = a^{mn}$)}
Simplify each expression. For Questions 15 and 17, also evaluate your answer as a single number.

\begin{qtwocol}
  \item $(3^2)^4$
  \qspace
  \item $(x^3)^5$
  \qspace
  \item $(2^2)^3$
  \qspace
  \item $(5^3)^2$
  \qspace
  \item $(2^4)^3$
  \qspace
  \item $(x^2)^6$
  \qspace
\end{qtwocol}

\subsection*{Writing exponentials with a common base}
\begin{qtwocol}
  \item Write 32 as a power of 2.
  \qspace
  \item Write 81 as a power of 3.
  \qspace
  \item Write 64 as a power of 2.
  \qspace
  \item Write $\dfrac{1}{9}$ as a power of 3.
  \qspace
  \item Rewrite $9^x$ as a power of 3 (in the form $3^{\,\square}$).
  \qspace
  \item Rewrite $\left(\dfrac{1}{3}\right)^x$ as a power of 3.
  \qspace
\end{qtwocol}

\subsection*{Combining related bases into a single base}
Some expressions mix two different bases that are both powers of the same prime (for example, $8 = 2^3$). Rewrite every term with that prime as the base, then use the multiplying or dividing law to combine them into a \textbf{single power}.

\begin{qtwocol}
  \item Simplify $2^{3x+1} \times 8^{2x+8}$ as a single power of 2.
  \qspace
  \item Simplify $9^{x} \times 3^{2x+5}$ as a single power of 3.
  \qspace
  \item Simplify $4^{x+2} \times 2^{3x-1}$ as a single power of 2.
  \qspace
  \item Simplify $\dfrac{16^{x+1}}{4^{x-2}}$ as a single power of 2.
  \qspace
  \item Simplify $27^{x} \times 3^{x+2}$ as a single power of 3.
  \qspace
  \item Simplify $\dfrac{8^{x+2}}{2^{x-1}}$ as a single power of 2.
  \qspace
\end{qtwocol}

\subsection*{Solving exponential equations}
Once both sides are written as a power of the same base, equate the exponents and solve.

\begin{qtwocol}
  \item Solve $2^x = 32$.
  \qspace
  \item Solve $5^x = \dfrac{1}{25}$.
  \qspace
  \item Solve $3^{2x-1} = 27$.
  \qspace
  \item Solve $4^{x} = 256$.
  \qspace
  \item Solve $2^{x-1} = \dfrac{1}{8}$.
  \qspace
  \item Solve $2^{3x+1} \times 8^{2x+8} = 2^{25}$. (Use your answer to Question 25.)
  \qspace
\end{qtwocol}

\newpage
\section{Problem Solving}

\noindent Many exponential equations are not immediately ready to solve. First \textbf{rearrange} the equation so the exponential term is on its own, then rewrite both sides with the same base, then equate exponents and solve.

\subsection*{Rearranging to a common base}
\begin{qtwocolstart}
  \item Solve $3^x + 5 = 14$.
  \qspace

  \item Solve $2 \times 5^x - 3 = 47$.
  \qspace

  \item Solve $\dfrac{4^x}{2} + 1 = 33$.
  \qspace

  \item Solve $50 - 2^x = -14$.
  \qspace

  \item Solve $9^x = 3^{x+4}$. (Both sides are already exponential --- rewrite the left-hand side using base 3 before equating exponents.)
  \qspace

  \item Solve $3^{x+3} \times 3^{2x} = 9^{x+4}$.
  \qspace
\end{qtwocolstart}

\subsection*{Leaving your answer in terms of $k$}
In these problems, $k$ is a fixed but unspecified number. Solve for $x$ as usual, leaving your final answer as an expression in terms of $k$.

\begin{qtwocol}
  \item Solve $3^{2x-1} = 3^{k+5}$ for $x$ in terms of $k$.
  \qspace

  \item Solve $4^x = 8^k$ for $x$ in terms of $k$. (Rewrite both sides using base 2 first.)
  \qspace

  \item Solve $2^{x} \times 4^{x} = 2^{k}$ for $x$ in terms of $k$.
  \qspace

  \item Solve $\dfrac{8^{x+1}}{2^{x-3}} = 2^{k}$ for $x$ in terms of $k$.
  \qspace

  \item Solve $9^{x} = 3^{k+4}$ for $x$ in terms of $k$.
  \qspace

  \item Solve $\dfrac{2^{x+5}}{4^{x-1}} = 2^{k}$ for $x$ in terms of $k$.
  \qspace
\end{qtwocol}

\newpage
\section{Enrichment}

\subsection*{Factoring out a separate prime}
When a base is a product of two primes (such as $6 = 2 \times 3$), dividing by a power of one of those primes can leave a power of the other. For example,
\[
  \frac{6^{x+3}}{3^{x+3}} = \left(\frac{6}{3}\right)^{x+3} = 2^{x+3}.
\]

\begin{qtwocolstart}
  \item Use this idea to solve: $\dfrac{6^{x+3}}{3^{x+3}} = 2^{5x+1}$.
  \qspace

  \item Solve: $\dfrac{10^{x+1}}{5^{x+1}} = 2^{3x-2}$.
  \qspace

  \item Solve: $\dfrac{15^{x}}{5^{x}} = 3^{2x-4}$.
  \qspace

  \item Solve: $\dfrac{14^{x-2}}{7^{x-2}} = 2^{3x-10}$.
  \qspace

  \item Solve: $\dfrac{21^{x}}{7^{x}} = 3^{2x-8}$.
  \qspace

  \item Solve: $\dfrac{33^{x-1}}{11^{x-1}} = 3^{2x+1}$.
  \qspace
\end{qtwocolstart}

\subsection*{All skills at once}

\begin{enumerate}[resume]
  \item Consider the equation
  \[
    \frac{(4^{x})^{2}}{2^{3}} \times \frac{20^{x+1}}{5^{x+1}} = 2^{2k-1}.
  \]
  \begin{enumerate}
    \item Simplify $\dfrac{(4^{x})^{2}}{2^{3}}$ as a single power of 2.
    \item Simplify $\dfrac{20^{x+1}}{5^{x+1}}$ as a single power of 2, by factoring out the prime 5.
    \item Combine your two results and equate exponents to solve for $x$ in terms of $k$.
    \item Check your formula by substituting $k = 6$ into your expression for $x$, then substituting both values back into the original equation to confirm both sides are equal.
  \end{enumerate}
  \qspace

  \item Consider the equation
  \[
    \frac{20^{x+1}}{5^{x+1}} \times 2^{3x} = 2^{k+2}.
  \]
  \begin{enumerate}
    \item Factor out the prime 5 to simplify $\dfrac{20^{x+1}}{5^{x+1}}$ as a single power of 2.
    \item Combine this with the $2^{3x}$ term, then equate exponents to solve for $x$ in terms of $k$.
    \item Check your formula by substituting $k = 10$ and confirming both sides of the original equation are equal.
  \end{enumerate}
  \qspace
\end{enumerate}

\subsection*{Challenge}
No scaffolding this time --- decide for yourself how to simplify this equation before you try to solve it.

\begin{enumerate}[resume]
  \item Solve for $x$: \quad $2^{x+3} \times 3^{x+1} = 2^{5} \times 3^{2x-1}$.
  \qspace
\end{enumerate}

\end{document}
